\documentclass[12pt, a4paper]{article}

\usepackage[T1]{fontenc}
\usepackage[utf8]{inputenc}
\usepackage{lmodern}
\usepackage[english]{babel}
\usepackage[left=2cm, right=2cm, top=2.5cm, bottom=2.5cm]{geometry}
\usepackage{booktabs}
\usepackage{longtable}
\usepackage{tabularx}
\usepackage{array}
\usepackage{multirow}
\usepackage{makecell}
\usepackage{xcolor}
\usepackage{colortbl}
\usepackage{hyperref}
\usepackage{setspace}
\usepackage[style=apa, backend=biber]{biblatex}
\addbibresource{references.bib}

\onehalfspacing

\definecolor{threadA}{RGB}{30, 90, 160}
\definecolor{threadB}{RGB}{180, 60, 30}
\definecolor{threadC}{RGB}{40, 130, 80}
\definecolor{threadD}{RGB}{130, 60, 150}
\definecolor{shadeA}{RGB}{230, 240, 255}
\definecolor{shadeB}{RGB}{255, 235, 225}
\definecolor{shadeC}{RGB}{225, 248, 230}
\definecolor{shadeD}{RGB}{245, 230, 255}

\hypersetup{colorlinks=true, citecolor=blue!60!black, linkcolor=blue!60!black}

\title{\textbf{Literature Overview Tables}\\
\large Time Perspective, Public Crisis, and Temporal Representation in LLMs}
\author{Lee Starrywalker}
\date{\today}

\begin{document}
\maketitle
\tableofcontents
\newpage

%% ================================================================
\section{Thread~I: Time Perspective Theory --- Core Literature}
%% ================================================================

Table~\ref{tab:tp_theory} presents the foundational and major empirical works on
Time Perspective (TP) theory, organized chronologically within sub-streams.

\begin{longtable}{@{}p{3.8cm} p{1.0cm} p{2.6cm} p{4.5cm} p{3.2cm}@{}}
\caption{Core Literature on Time Perspective Theory}
\label{tab:tp_theory}\\
\toprule
\textbf{Author(s)} & \textbf{Year} & \textbf{Venue / IF} &
\textbf{Core Contribution} & \textbf{Key Findings / Methods} \\
\midrule
\endfirsthead
\multicolumn{5}{l}{\small\emph{Table~\ref{tab:tp_theory} continued\ldots}}\\
\toprule
\textbf{Author(s)} & \textbf{Year} & \textbf{Venue / IF} &
\textbf{Core Contribution} & \textbf{Key Findings / Methods} \\
\midrule
\endhead
\midrule
\multicolumn{5}{r}{\small\emph{Continued on next page}}\\
\endfoot
\bottomrule
\endlastfoot

%% ─── Sub-stream 1a: Theoretical Precursors ───────────────────────────────
\rowcolor{shadeA}
\multicolumn{5}{l}{\textcolor{threadA}{\textbf{1a. Theoretical Precursors}}}\\
\midrule

\textcite{lewin1951field}
  & 1951 & Monograph
  & Psychological field theory; introduced \emph{life space} with temporal dimension
  & Conceptual; behavior shaped by past representations and future anticipations \\[4pt]

\textcite{nuttin1985future}
  & 1985 & Monograph
  & Defined \emph{temporal range} and \emph{temporal depth} within motivational theory
  & Theoretical; capacity to concretize future goals drives self-regulation \\[6pt]

%% ─── Sub-stream 1b: Instrument Development ──────────────────────────────
\rowcolor{shadeA}
\multicolumn{5}{l}{\textcolor{threadA}{\textbf{1b. Instrument Development \& Validation}}}\\
\midrule

\textcite{strathman1994cfc}
  & 1994 & \emph{JPSP} (IF\,8.0)
  & Developed 12-item Consideration of Future Consequences (CFC) scale
  & 7 samples; predicts health behavior, env.\ protection, risk decisions \\[4pt]

\textcite{zimbardo1999putting}
  & 1999 & \emph{JPSP} (IF\,8.0)
  & Proposed ZTPI five-factor structure: PN, PP, PH, PF, F
  & 7 independent samples; validated across 24 countries \\[4pt]

\textcite{worrell2016theoretical}
  & 2016 & \emph{EJPA}
  & Proposed theory-driven ZTPI-TP retaining temporally referential items only
  & 4 countries, 5 samples; superior structural validity; PF/PH issues remain \\[4pt]

\textcite{temple2019zimbardo}
  & 2019 & \emph{Time \& Society}
  & Tested ZTPI-20/-17/-15 short versions across national samples
  & 5 samples; all short versions show factor validity problems; theory-driven revision recommended \\[4pt]

\textcite{wiberg2021systematic}
  & 2021 & \emph{Front.\ Psychol.}
  & IRT-based systematic review to identify robust ZTPI items
  & Identified most psychometrically stable items per subscale \\[6pt]

%% ─── Sub-stream 1c: Meta-Analytic Synthesis ─────────────────────────────
\rowcolor{shadeA}
\multicolumn{5}{l}{\textcolor{threadA}{\textbf{1c. Meta-Analytic Synthesis \& Self-Regulation}}}\\
\midrule

\textcite{boniwell2004balancing}
  & 2004 & Book Chapter
  & Proposed Balanced Time Perspective (BTP) as normative target
  & Conceptual; balance among PP, moderate PH, and F as optimal profile \\[4pt]

\textcite{hall2007temporal}
  & 2007 & \emph{Health Psychol.\ Rev.}
  & Temporal Self-Regulation Theory (TST): TP within goal setting, monitoring, pursuit
  & Theoretical model; formalized three-stage self-regulation pathway \\[4pt]

\textcite{kooij2018future}
  & 2018 & \emph{J.\ Appl.\ Psychol.} (IF\,9.4)
  & Systematic review of Future TP across applied domains
  & $k=212$ studies; FTP predicts achievement, health, retirement after controlling Big-Five \\[4pt]

\textcite{baird2021psychological}
  & 2021 & \emph{Psychol.\ Bull.} (IF\,20.85)
  & Meta-analysis of TP's self-regulatory effects
  & $k=378$, 2000 ES, RVE; $r_+(\text{goal setting})=0.25$; mediators: goal monitoring, pursuit, ability \\[4pt]

\textcite{transdiagnostic2024}
  & 2024 & \emph{Int'l J.\ Psychol.}
  & Transdiagnostic meta-analysis of ZTPI dimensions
  & 132 articles, 367 ES (1990--2023); Future TP as protective factor across internalizing, ADHD, substance use \\[6pt]

%% ─── Sub-stream 1d: Socioemotional Selectivity ──────────────────────────
\rowcolor{shadeA}
\multicolumn{5}{l}{\textcolor{threadA}{\textbf{1d. Socioemotional Selectivity Theory (SST)}}}\\
\midrule

\textcite{carstensen1999taking}
  & 1999 & \emph{Am.\ Psychologist}
  & SST: perceived remaining time shapes goal priority (knowledge vs.\ emotion)
  & SARS case: young adults shift TP under crisis-induced temporal finitude \\[4pt]

\textcite{carstensen2006influence}
  & 2006 & \emph{Science} (IF\,56.9)
  & Extended SST across lifespan and situational time constraints
  & Crisis triggers TP shifts; even non-elderly shift to emotionally meaningful goals \\[6pt]

%% ─── Sub-stream 1e: Clinical Applications ───────────────────────────────
\rowcolor{shadeA}
\multicolumn{5}{l}{\textcolor{threadA}{\textbf{1e. Clinical \& Crisis Applications}}}\\
\midrule

\textcite{henson2006time}
  & 2006 & \emph{J.\ Behav.\ Med.}
  & TP predicts health behaviors across multiple domains
  & High Future $\to$ positive health behavior; high PH $\to$ risk behavior (cross-domain) \\[4pt]

\textcite{haas2023time}
  & 2023 & \emph{IJERPH}
  & ZTPI + CFC jointly predict COVID-19 vaccination intention
  & $n=232$; PP and BTP deviation \emph{negatively} predict intent (counter-intuitive) \\[4pt]

\textcite{sword2013time}
  & 2014 & \emph{J.\ Loss Trauma}
  & Time Perspective Therapy (TPT) for PTSD
  & 29 veterans, 4 years: 89\% reduced depression, 70\% anxiety, 52\% trauma \\[4pt]

\textcite{yu2025ztpt}
  & 2025 & \emph{J.\ Bengbu Med.\ Univ.}
  & ZTPT + mindfulness RCT for PTSD
  & $n=80$; reduced PCL-C, PSQI, HAMA, HAMD-17; modulates Cor, NE, 5-HT, IL-6 \\

\end{longtable}

\newpage

%% ================================================================
\section{Thread~II: LLMs in Public Crisis Response}
%% ================================================================

Table~\ref{tab:llm_crisis} documents the literature on LLM applications and limitations
in public crisis contexts.

\begin{longtable}{@{}p{3.8cm} p{1.0cm} p{2.6cm} p{4.5cm} p{3.2cm}@{}}
\caption{LLMs in Public Crisis Response: Applications, Benchmarks, and Limitations}
\label{tab:llm_crisis}\\
\toprule
\textbf{Author(s)} & \textbf{Year} & \textbf{Venue} &
\textbf{Core Contribution} & \textbf{Key Findings / Methods} \\
\midrule
\endfirsthead
\multicolumn{5}{l}{\small\emph{Table~\ref{tab:llm_crisis} continued\ldots}}\\
\toprule
\textbf{Author(s)} & \textbf{Year} & \textbf{Venue} &
\textbf{Core Contribution} & \textbf{Key Findings / Methods} \\
\midrule
\endhead
\midrule
\multicolumn{5}{r}{\small\emph{Continued on next page}}\\
\endfoot
\bottomrule
\endlastfoot

%% ─── Sub-stream 2a: Applications ────────────────────────────────────────
\rowcolor{shadeB}
\multicolumn{5}{l}{\textcolor{threadB}{\textbf{2a. Crisis Applications of LLMs}}}\\
\midrule

\textcite{geoknowledge2023}
  & 2023 & EMNLP Industry
  & Geo-knowledge-guided GPT for disaster tweet location extraction
  & 22-shot; outperforms 9 baselines on Hurricane Harvey data \\[4pt]

\textcite{floodbrain2023}
  & 2023 & \emph{IJDRR}
  & FloodBrain: RAG-LLM automated flood-impact report generation
  & Human-rated quality comparable to expert reports; reduces coordination time \\[4pt]

\textcite{aiinfodemic2023}
  & 2023 & \emph{J.\ Med.\ Lib.\ Assoc.}
  & LLMs as potential infodemic amplifiers in COVID-19 context
  & Review; hallucination risk in public health communication documented \\[6pt]

%% ─── Sub-stream 2b: Benchmarks ──────────────────────────────────────────
\rowcolor{shadeB}
\multicolumn{5}{l}{\textcolor{threadB}{\textbf{2b. Crisis Reasoning Benchmarks}}}\\
\midrule

\textcite{diallo2025response}
  & 2025 & ACL 2025
  & RESPONSE benchmark: LLM disaster common-sense reasoning
  & 1,789 instances, 6,037 QA groups; GPT-4 only 37\% accuracy on immediate frame \\[4pt]

\textcite{wei2025time}
  & 2025 & NeurIPS 2025 \textit{Spotlight}
  & TIME benchmark: multi-level temporal reasoning in real-world scenarios
  & 38,522 QA pairs, 11 subtasks; temporal retrieval is core bottleneck \\[4pt]

\textcite{tdbench2026}
  & 2026 & ICLR 2026
  & TDBench: temporal database reasoning for LLMs
  & $\approx$21.7\% reasoning hallucinations (correct answer, wrong time evidence) \\[6pt]

%% ─── Sub-stream 2c: Temporal Limitations ───────────────────────────────
\rowcolor{shadeB}
\multicolumn{5}{l}{\textcolor{threadB}{\textbf{2c. Temporal Limitations of LLMs}}}\\
\midrule

\textcite{qiu2024temporally}
  & 2024 & NAACL 2024
  & Systematic probe of temporal grounding in LLMs
  & $\geq$27.23\% self-consistency failures; scaling does not improve; sentence order $\neq$ event order \\[4pt]

\textcite{ticktoc2025}
  & 2025 & EMNLP 2025
  & TicToc: reveals ``temporal blindness'' in LLM agents
  & 18 LLMs tested; best alignment 65\% with timestamps; $<$4\% reasoning traces mention time \\[4pt]

\textcite{cheng2024dated}
  & 2024 & COLM 2024 \textit{(Outstanding)}
  & Concept of effective cutoff vs.\ reported cutoff
  & WIKISPAN/NEWSSPAN probing; effective cutoff precedes official by months--years \\[4pt]

\textcite{zhu2025freshbench}
  & 2025 & NAACL 2025
  & FreshBench: temporal generalization evaluation axis
  & LLM performance declines systematically over time; significant past-bias demonstrated \\[6pt]

%% ─── Sub-stream 2d: Cognitive Bias Inheritance ──────────────────────────
\rowcolor{shadeB}
\multicolumn{5}{l}{\textcolor{threadB}{\textbf{2d. Cognitive Bias \& Early Cognitive Research}}}\\
\midrule

\textcite{echterhoff2024cognitive}
  & 2024 & EMNLP Findings 2024
  & BiasBuster framework for cognitive bias in LLM decision-making
  & 13,465 prompts, 13 bias types; anchoring and framing effects confirmed across models \\[4pt]

\textcite{testoftimereasoning2025}
  & 2025 & Google Research
  & Systematic decomposition of LLM temporal reasoning factors
  & Graph structure, question type, fact ordering all modulate temporal performance \\[4pt]

\textcite{bogdanov2026multiturn}
  & 2026 & ICLR 2026 Workshop
  & Survival-analysis of multi-turn LLM temporal coherence degradation
  & Early failures: impulsivity-driven; later failures: fatigue/cost-benefit framing \\

\end{longtable}

\newpage

%% ================================================================
\section{Thread~III: Temporal Representation in NLP --- Historical Trajectory}
%% ================================================================

Table~\ref{tab:tr_history} traces the development of temporal representation in NLP from
annotation standards through the neural era.

\begin{longtable}{@{}p{3.8cm} p{1.0cm} p{2.6cm} p{4.5cm} p{3.2cm}@{}}
\caption{Temporal Representation in NLP: Historical Development}
\label{tab:tr_history}\\
\toprule
\textbf{Author(s)} & \textbf{Year} & \textbf{Venue} &
\textbf{Core Contribution} & \textbf{Key Findings / Methods} \\
\midrule
\endfirsthead
\multicolumn{5}{l}{\small\emph{Table~\ref{tab:tr_history} continued\ldots}}\\
\toprule
\textbf{Author(s)} & \textbf{Year} & \textbf{Venue} &
\textbf{Core Contribution} & \textbf{Key Findings / Methods} \\
\midrule
\endhead
\midrule
\multicolumn{5}{r}{\small\emph{Continued on next page}}\\
\endfoot
\bottomrule
\endlastfoot

%% ─── Sub-stream 3a: Annotation Standards ───────────────────────────────
\rowcolor{shadeC}
\multicolumn{5}{l}{\textcolor{threadC}{\textbf{3a. Annotation Standards: TimeML \& TimeBank}}}\\
\midrule

\textcite{pustejovsky2003timeml}
  & 2003 & IWCS-5
  & TimeML: specification for TIMEX3, EVENT, TLINK, SLINK
  & Canonical standard for $>$10 years of temporal IE research; TimeBank corpus \\[4pt]

\textcite{pustejovsky2010iso}
  & 2010 & LAW-IV
  & ISO-TimeML: internationalized TimeML specification; added RLINK
  & Extended standard; enables cross-lingual temporal annotation \\[4pt]

\textcite{llorens2010timebank}
  & 2010 & LREC 2010
  & TRIOS-TimeBank: added missing events, new relation categories
  & Improved annotation coverage and consistency over original TimeBank \\[4pt]

\textcite{minard2016meantime}
  & 2016 & LREC 2016
  & MEANTIME: multilingual event and time corpus; added ELINK
  & Extends temporal structure to multilingual news; richer semantic annotation \\[6pt]

%% ─── Sub-stream 3b: TempEval Benchmarks ────────────────────────────────
\rowcolor{shadeC}
\multicolumn{5}{l}{\textcolor{threadC}{\textbf{3b. TempEval Shared Tasks (2007--2016)}}}\\
\midrule

\textcite{pustejovsky2003timeml}
  & 2007 & SemEval-2007
  & TempEval-1: first evaluation framework for TIMEX3 and event-time relations
  & Established shared evaluation; early ML systems benchmarked \\[4pt]

\textcite{tempeval2010}
  & 2010 & SemEval-2010
  & TempEval-2: extended to 6 languages
  & Multilingual generalizability confirmed; cross-lingual temporal models evaluated \\[4pt]

\textcite{heideltime2013}
  & 2013 & SemEval-2013
  & HeidelTime: best TIMEX3 extraction/normalization in TempEval-3
  & Rule-based; best on English and Spanish; multilingual coverage \\[4pt]

\textcite{sutime2012}
  & 2012 & LREC 2012
  & SuTime (Stanford CoreNLP): deterministic temporal tagger
  & Surpassed contemporaneous systems on TempEval-2; rule-based normalization \\[4pt]

\textcite{tempeval2015}
  & 2015 & SemEval-2015
  & QA TempEval: extrinsic QA-based evaluation paradigm shift
  & Best system recall $\approx$30\%; exposed gap between intrinsic and extrinsic performance \\[4pt]

\textcite{tempeval2016}
  & 2016 & SemEval-2016
  & Clinical TempEval: clinical text temporal extraction
  & Near-human performance achieved; domain-transfer gap identified \\[6pt]

%% ─── Sub-stream 3c: Neural / Pre-trained Models ─────────────────────────
\rowcolor{shadeC}
\multicolumn{5}{l}{\textcolor{threadC}{\textbf{3c. Neural Era: Datasets and Pre-trained Models}}}\\
\midrule

\textcite{ning2018matres}
  & 2018 & ACL 2018
  & MATRES: multi-axis event start-point ordering annotation scheme
  & Cleaner annotation; multi-axis reduces label ambiguity \\[4pt]

\textcite{ning2019improved}
  & 2019 & EMNLP 2019
  & ELMo/BERT + Siamese KB encoder + ILP for temporal relation extraction
  & $\approx$10 pp gain; 25\% error reduction over SVM baselines \\[4pt]

\textcite{torque2020}
  & 2020 & EMNLP 2020
  & TORQUE: SQuAD-style temporal ordering reading comprehension
  & Bridges classical IE and QA paradigms; human--model gap highlighted \\[4pt]

\textcite{tacolm2020}
  & 2020 & ACL 2020
  & TacoLM: temporal commonsense pre-training (explicit + implicit mining)
  & Outperforms BERT on TimeBank and HiEVE; temporal commonsense as pre-training signal \\[4pt]

\textcite{timedial2021}
  & 2021 & ACL 2021
  & TimeDial: dialogic temporal commonsense reasoning benchmark
  & Exposes turn-level temporal tracking deficits in PLMs \\[4pt]

\textcite{bitembert2024}
  & 2024 & CIKM 2024
  & BiTimeBERT 2.0: ETAMLM + date prediction + entity replacement pre-training
  & Significant gains; training cost reduced by $\approx$53\% \\[4pt]

\textcite{multitask2023}
  & 2023 & \emph{Sci.\ Rep.}
  & Multi-task learning with temporal relation extraction auxiliary objective
  & +6.4 pp Korean, +3.6 pp English NLU accuracy \\

\end{longtable}

\newpage

%% ================================================================
\section{Thread~IV: Temporal Representation in LLMs --- Six Pathways \& Obstacles}
%% ================================================================

Table~\ref{tab:tr_llm} covers recent work on TR within large language models, organized
by research pathway and the dual obstacles to direct study.

\begin{longtable}{@{}p{3.8cm} p{1.0cm} p{2.6cm} p{4.5cm} p{3.2cm}@{}}
\caption{Temporal Representation in Large Language Models: Pathways and Obstacles}
\label{tab:tr_llm}\\
\toprule
\textbf{Author(s)} & \textbf{Year} & \textbf{Venue} &
\textbf{Core Contribution} & \textbf{Key Findings / Methods} \\
\midrule
\endfirsthead
\multicolumn{5}{l}{\small\emph{Table~\ref{tab:tr_llm} continued\ldots}}\\
\toprule
\textbf{Author(s)} & \textbf{Year} & \textbf{Venue} &
\textbf{Core Contribution} & \textbf{Key Findings / Methods} \\
\midrule
\endhead
\midrule
\multicolumn{5}{r}{\small\emph{Continued on next page}}\\
\endfoot
\bottomrule
\endlastfoot

%% ─── Pathway (a): Representational Probing ──────────────────────────────
\rowcolor{shadeD}
\multicolumn{5}{l}{\textcolor{threadD}{\textbf{Pathway (a): Representational Probing}}}\\
\midrule

\textcite{gurnee2024language}
  & 2024 & ICLR 2024
  & Probed space--time representations in Llama-2 (7B/13B/70B)
  & Approximately linear temporal encoding in middle layers; ``time neurons'' causally validated \\[4pt]

\textcite{song2025temporal}
  & 2025 & NeurIPS 2025
  & SAE + identifiable temporal--instantaneous causal representation learning framework
  & Theoretical identifiability for LLM concept space; meaningful temporal concepts uncovered \\[4pt]

\textcite{thukral2021probing}
  & 2021 & BlackboxNLP 2021
  & NLI-based temporal ordering probes for RoBERTa/DeBERTa
  & Direct accuracy only 52--56\%; models rely on distribution, not intrinsic temporal knowledge \\[6pt]

%% ─── Pathway (b): Time-Aware Pretraining ────────────────────────────────
\rowcolor{shadeD}
\multicolumn{5}{l}{\textcolor{threadD}{\textbf{Pathway (b): Time-Aware Pretraining}}}\\
\midrule

\textcite{dhingra2022talm}
  & 2022 & TACL
  & TALM: joint modeling of text and timestamps
  & Improved seen-fact memorization and future-period calibration; efficient model refreshing \\[6pt]

%% ─── Pathway (c): Temporal Knowledge Editing ───────────────────────────
\rowcolor{shadeD}
\multicolumn{5}{l}{\textcolor{threadD}{\textbf{Pathway (c): Temporal Knowledge Editing}}}\\
\midrule

\textcite{yin2024history}
  & 2024 & AAAI 2024
  & TKE task + AToKe benchmark; METO framework
  & Existing editors cause catastrophic forgetting; METO jointly preserves historical and new knowledge \\[6pt]

%% ─── Pathway (d): Reasoning Enhancement ────────────────────────────────
\rowcolor{shadeD}
\multicolumn{5}{l}{\textcolor{threadD}{\textbf{Pathway (d): Temporal Reasoning Enhancement}}}\\
\midrule

\textcite{tiser2025}
  & 2025 & ACL 2025
  & TISER: four-stage timeline self-reflection pipeline
  & Small fine-tuned models $>$ larger closed models; test-time compute scaling effective \\[4pt]

\textcite{zhang2024not}
  & 2024 & EMNLP Findings 2024
  & Narrative-of-Thought (NoT): temporal reasoning as graph generation
  & F1 improvement 16--71\% on TB-Dense; $<$10B models reach GPT-3.5 performance \\[4pt]

\textcite{huang2023etre}
  & 2023 & ACL 2023
  & ETRE: unified framework via temporal-point logical expressions
  & Outperforms strong baselines on TB-Dense and MATRES \\[6pt]

%% ─── Pathway (e): Temporal Bias Auditing ────────────────────────────────
\rowcolor{shadeD}
\multicolumn{5}{l}{\textcolor{threadD}{\textbf{Pathway (e): Temporal Bias Auditing}}}\\
\midrule

\textcite{setclock2024}
  & 2024 & ACL Findings 2024
  & Training-time bias in PLMs; local temporal alignment correction
  & Models favor temporally dense training windows; alignment fine-tuning improves QA \\[4pt]

\textcite{cheng2024dated}
  & 2024 & COLM 2024
  & Effective cutoff concept; WIKISPAN/NEWSSPAN probing methodology
  & Effective cutoffs precede official cutoffs; CommonCrawl staleness + dedup failures as root cause \\[6pt]

%% ─── Pathway (f): Temporal Generalization ───────────────────────────────
\rowcolor{shadeD}
\multicolumn{5}{l}{\textcolor{threadD}{\textbf{Pathway (f): Temporal Generalization \& Continual Learning}}}\\
\midrule

\textcite{zhu2025freshbench}
  & 2025 & NAACL 2025
  & FreshBench: dynamic benchmark from real-world predictions
  & Systematic performance decline over time; significant past-bias \\[4pt]

\textcite{bogdanov2026multiturn}
  & 2026 & ICLR Workshop
  & Survival analysis of LLM temporal coherence across turns
  & Coherence degrades with turns; early = impulsive, late = fatigue/cost-benefit \\[6pt]

%% ─── Evaluation Benchmarks ──────────────────────────────────────────────
\rowcolor{shadeD}
\multicolumn{5}{l}{\textcolor{threadD}{\textbf{Evaluation Benchmarks (TR in LLMs)}}}\\
\midrule

\textcite{timebench2024}
  & 2024 & ACL 2024
  & TimeBench: comprehensive temporal reasoning evaluation suite
  & Multi-type temporal tasks; reveals different failure modes per task type \\[4pt]

\textcite{tram2024}
  & 2024 & ACL Findings 2024
  & TRAM: temporal reasoning ability benchmarking for LLMs
  & Fine-grained temporal reasoning taxonomy; LLMs weakest on temporal arithmetic \\

\end{longtable}

\newpage

%% ================================================================
\section{Cross-Thread Synthesis: Key Gaps and Proposed Pathways}
%% ================================================================

Table~\ref{tab:gaps} provides a synthesized view of the critical research gaps identified
across the four literature threads, serving as the structural motivation for the proposed
indirect TR study strategy.

\begin{longtable}{@{}p{3.5cm} p{3.5cm} p{3.5cm} p{4.2cm}@{}}
\caption{Key Research Gaps at the Intersection of the Three Threads}
\label{tab:gaps}\\
\toprule
\textbf{Gap Dimension} & \textbf{Description} &
\textbf{Evidence from Literature} & \textbf{Proposed Remedy} \\
\midrule
\endfirsthead
\toprule
\textbf{Gap Dimension} & \textbf{Description} &
\textbf{Evidence from Literature} & \textbf{Proposed Remedy} \\
\midrule
\endhead
\bottomrule
\endlastfoot

TP--LLM Measurement Gap
  & ZTPI/CFC not applied to LLM capability assessment
  & \textcite{zimbardo1999putting}; \textcite{tram2024}
  & ZTPI-aligned LLM evaluation dataset for crisis scenarios \\[6pt]

Direct TR Obstacles
  & Black-box APIs + white-box complexity prevent direct probing
  & \textcite{ticktoc2025}; \textcite{tdbench2026}; \textcite{song2025temporal}
  & Indirect study via TP-grounded behavioral evaluation \\[6pt]

Temporal Grounding Deficit
  & $\geq$27.23\% self-consistency failures; sentence $\neq$ event order
  & \textcite{qiu2024temporally}
  & Temporal self-consistency testing protocol \\[6pt]

Crisis Temporal Cognition Gap
  & TP-dimension biases in LLM crisis outputs unstudied
  & \textcite{diallo2025response}; \textcite{wei2025time}
  & TP prompt injection + RESPONSE-aligned evaluation \\[6pt]

Effective Cutoff Distortion
  & Official cutoffs precede effective cutoffs by years
  & \textcite{cheng2024dated}; \textcite{zhu2025freshbench}
  & Temporally aware RAG with cutoff-corrected retrieval \\[6pt]

Human--LLM TP Asymmetry
  & Humans shift TP under crisis; LLMs have no such mechanism
  & \textcite{carstensen1999taking}; \textcite{ticktoc2025}
  & Simulate SST-like context shifts via structured prompting \\

\end{longtable}

\printbibliography

\end{document}
